\section{Related Work}
\label{sec:related}

\paragraph{LLM serving systems.}
Orca introduced iteration-level scheduling~\cite{yu2022orca}; \vllm added paged KV-cache
management~\cite{kwon2023vllm}; \sglang added prefix sharing through a radix-tree cache and a
frontend language for structured programs~\cite{zheng2024sglang}; \trt provides NVIDIA's
optimised kernels and in-flight batching~\cite{trtllm}. Research systems improve on specific
aspects: chunked prefill to bound decode stalls~\cite{agrawal2024sarathi}, prefill--decode
disaggregation~\cite{zhong2024distserve,patel2024splitwise}, preemptive and fair
scheduling~\cite{wu2023fastserve}, live migration across instances~\cite{sun2024llumnix}, and
KV-cache-centric architectures~\cite{qin2025mooncake}. These papers evaluate their own design
against baselines; we instead compare widely deployed engines as black boxes, through the
interface applications use. A survey of serving techniques is given by Miao
et~al.~\cite{miao2023survey}.

\paragraph{Benchmarking and evaluating LLM inference.}
MLPerf Inference~\cite{reddi2020mlperf} defines scenarios and latency constraints for inference
benchmarks but targets submitted, tuned systems rather than default deployments under realistic
mixed traffic. Etalon~\cite{agrawal2024etalon} argues that TTFT, TBT and TPOT individually miss
user-visible stalls and proposes a fluidity metric; Andes~\cite{liu2024andes} defines quality of
experience for streamed text. Vidur~\cite{agrawal2024vidur} simulates serving configurations to
search for good deployments. Tools such as LLMPerf~\cite{llmperf} measure latency and throughput
of serving endpoints. Our study complements these with a controlled cross-engine comparison under
saturation, an explanation of the differences in terms of admission and eviction behaviour, and a
set of methodology pitfalls observed in practice.

\paragraph{Performance-engineering methodology.}
Open- and closed-loop load generators produce different conclusions about the same
system~\cite{schroeder2006open}, and latency recorders that measure from dispatch suffer from
coordinated omission~\cite{tene2015latency}. Little's law~\cite{little1961} relates queue length,
arrival rate and residence time; we use it to explain when the scheduler's sequence limit becomes
the bottleneck. Our methodology lessons (\cref{sec:lessons}) apply these principles to LLM serving,
where long, variable-length generations make them especially consequential.

\todo{add recent empirical comparisons of serving engines (2025--2026) and position against them}
